% !PRINTED 325
%<<CONT
and represent the value $X:Y$ by a variable point of this circle,
and denote the points in question by $\alpha, x, X$, then
$\Phi$ passes through zero when $X$ goes from $x$ forwards or
backwards to the nearest points, which we shall denote by $\alpha_1$ and
$\alpha_2$. If then $X_1$, $X_2$ are the roots of
equation (9), we have the following arrangement on the circle

\weq{10}{\alpha_1,\ X_1,\ x,\ X_2,\ \alpha_2,}

\gesperrt{and $(X_1, X_2)$ represents an interval in which the arbitrarily
chosen point $x$ lies, but no root point of the equation $f = 0$.
$X_1$ and $X_2$ are therefore nearer to the two nearest roots
$\alpha_1, \alpha_2$ than the initial value $x$.}

This holds whatever position the second arbitrary point $\xi$
may have. The points $X_1, X_2$ will indeed depend
on it; but they will always lie in the interval $(\alpha_1, x)$ and $(x, \alpha_2)$.

There now arises the question: how is the point $\xi$ to be chosen so that
$X_1$ may lie as near as possible to $\alpha_1$, or $X_2$ as near as possible to $\alpha_2$;
or rather, since the knowledge of $\xi$ itself does not matter,
which is the point $X_1$ lying as near as possible to $\alpha_1$, and the
point $X_2$ lying as near as possible to $\alpha_2$?

If we think of the point $X$ as given, then (9) is a
quadratic equation for the point $\xi$; to every position of $X$
there thus belong two positions of $\xi$. But only such values of
$X_1$ or $X_2$ are admissible in (10) for which this quadratic
equation yields real roots. Hence only such values of
$X_1$ and $X_2$ can occur for which the discriminant of
(9) with respect to $\xi$ is positive, and we obtain the limiting positions
of $X_1$ and $X_2$ when we set the discriminant of (9) equal to
zero. But the coefficients of $\xi^2$, $2\xi\eta$, $\eta^2$ in (9) are
\[
\frac{f'(x)^2 - y^2 H}{nf^2}(Xy - Yx)^2 - Y^2
\]
\[
\frac{f'(x)f'(y) + xyH}{nf^2}(Xy - Yx)^2 + XY
\]
\[
\frac{f'(y)^2 - x^2 H}{nf^2}(Xy - Yx)^2 - X^2,
\]
and the discriminant is obtained when from the square
%CONT>>
